\begin{tikzfig}
\begin{adjustbox}{max width=\linewidth}\begin{tikzpicture}[fig]
% Базовый
\node[font=\sffamily\small\bfseries, anchor=west] at (0,1.3) {Пионер Базовый};
\node[forange, text width=44mm, minimum height=15mm] (pc1) at (2.7,0) {\textbf{Ваш компьютер}\\ Python, \texttt{pioneer\_sdk}\\ Pioneer Station};
\node[forange, text width=44mm, minimum height=15mm] (ap1) at (11.2,0) {\textbf{Автопилот}\\ Lua-скрипт, загруженный на борт};
\draw[farrow] (pc1) -- node[flab, above=1pt] {Wi-Fi или USB} (ap1);
% Мини 2
\node[font=\sffamily\small\bfseries, anchor=west] at (0,-2.3) {Пионер Мини~2};
\node[fgrey, text width=36mm, minimum height=15mm] (pc2) at (2.0,-4) {\textbf{Ваш компьютер}\\ браузер: Pioneer Code, Bricks};
\node[forange, text width=44mm, minimum height=15mm] (bc) at (8.7,-4) {\textbf{Бортовой компьютер}\\ Python, \texttt{pioneer\_sdk2}, \texttt{pioneer\_rknn}; блоки Bricks};
\node[fgrey, text width=22mm, minimum height=15mm] (ap2) at (13.6,-4) {\textbf{Автопилот}};
\draw[farrow] (pc2) -- node[flab, above=1pt, align=center, font=\sffamily\footnotesize] {Wi-Fi \\ дрона} (bc);
\draw[farrow] (bc) -- (ap2);
\begin{scope}[on background layer]
  \node[draw=Accent, dashed, rounded corners=3pt, line width=0.8pt, fit=(bc) (ap2), inner sep=6pt, label={[font=\sffamily\footnotesize, text=Accent, anchor=south east]north east:на борту дрона}] {};
  \node[draw=Accent, dashed, rounded corners=3pt, line width=0.8pt, fit=(ap1), inner sep=6pt, label={[font=\sffamily\footnotesize, text=Accent, anchor=south east]north east:на борту дрона}] {};
\end{scope}
\end{tikzpicture}\end{adjustbox}
\captionof{figure}{Где выполняется программа. У Базового Lua работает внутри автопилота, а Python --- на вашем компьютере. У Мини~2 Python работает на самом дроне, а компьютер нужен лишь для записи и запуска. Оранжевая заливка --- место, где выполняется ваш код. Симулятор воспроизводит те же программные интерфейсы вместо дрона.}\label{fig:where-code}
\end{tikzfig}
